% 摘要
\begin{abstract}
\noindent
政企场景下的大模型智能体被普遍接入工具调用能力,常见攻击方式从显式注入转向业务话术诱导下的过度代理。针对该威胁,本报告给出三层混合防护体系。第一层为结构化规则引擎,覆盖外发邮箱白名单、写入路径白名单、SQL 写操作、定时任务、审计日志清理等 14 类高风险动作。第二层为结构化 LLM 评审,在注入识别基础上引入过度代理检测与意图对齐判断两个维度,工程上采用异常即阻断的封闭失败模式。第三层为真实环境接入验证,在真实 OpenClaw 类智能体上完成端到端跑通。实验在 66 条真实业务话术诱导样本上,攻击拦截率从 86.4\% 提升至 100\%,在 25 条正常业务白样本上误报率为 0\%。

\vspace{1em}
\noindent\textbf{关键词:} 大模型智能体;提示词注入;工具调用安全;过度代理;纵深防御
\end{abstract}

\newpage
